\documentclass[11pt]{article}
\usepackage[T1]{fontenc}
\usepackage[utf8]{inputenc}
\usepackage{lmodern}
\usepackage[margin=1in]{geometry}
\usepackage{amsmath,amssymb,amsthm}
\usepackage{booktabs,longtable,array}
\usepackage{xcolor,listings}
\usepackage{microtype}
\usepackage{etoolbox,needspace}
\usepackage{xurl}
\usepackage{hyperref}
\setlength{\emergencystretch}{2em}
\setcounter{tocdepth}{1}
\hypersetup{colorlinks=true,linkcolor=blue!45!black,urlcolor=blue!45!black,
  citecolor=blue!45!black,pdftitle={Rank-N and Impredicative Code Synthesis in Djex and Leant}}
\lstset{basicstyle=\ttfamily\small,columns=fullflexible,keepspaces=true,
  breaklines=true,breakatwhitespace=false,showstringspaces=false,
  frame=single,rulecolor=\color{black!20},xleftmargin=0.5em,xrightmargin=0.5em,
  aboveskip=0.8em,belowskip=0.8em}
\BeforeBeginEnvironment{lstlisting}{\par\noindent\begin{minipage}{\linewidth}}
\AfterEndEnvironment{lstlisting}{\end{minipage}\par}
\pretocmd{\subsection}{\Needspace{6\baselineskip}}{}{}
\pretocmd{\section}{\Needspace{10\baselineskip}}{}{}
\newcommand{\code}[1]{\texttt{#1}}
\newcommand{\FV}{\operatorname{FV}}
\newcommand{\ty}{\mathsf{Type}}
\newcommand{\ctx}{\Gamma}
\newtheorem{proposition}{Proposition}
\title{Rank-N and Impredicative Code Synthesis\\in Djex and Leant}
\author{Djex and Leant contributors}
\date{September 2026}

\begin{document}
\maketitle

\begin{abstract}
This document describes the shared higher-rank synthesis architecture, the
structural instantiation and unification improvements, the Haskell and Lean
checking boundaries, and the acceptance corpus derived from
\path{docs/examples/Church.hs}. The corpus contains 346 top-level signatures
and four local signatures, including 51 partial source signatures whose holes
are resolved by GHC. Every case is tested without supplying a reference
implementation. There are 315 cases with no named value providers, 16 cases
requiring a concrete integer constant, and 19 inherently partial cases tested
under an explicit element-default assumption. Both Djinn and Exference produce
350 candidates, and GHC checks all 700 generated implementations at both their
expanded query types and their original alias-preserving acceptance types.
The 19 partial Haskell cases per engine additionally receive an explicitly
partial wrapper at the exact original signature; bottom is inserted only at
the documented default argument after synthesis.
Independent live Lean runs on one unchanged executable also produce all 350
universe-correct counterparts per engine. All 700 exact displayed terms pass
standalone kernel replay with empty axiom inventories. Lean retains the
explicit default premise for the 19 partial source cases. These results establish a broad, reproducible
practical acceptance boundary; they do not assert a terminating decision
procedure for arbitrary System~F inhabitation or semantic equivalence with the
reference functions.
\end{abstract}

\newpage
\tableofcontents
\newpage

\section{The problem and the supported contract}

A synthesis request consists of a type $T$, a context $\ctx$ of available
values, and operational search limits. A successful result is a term $e$ such
that
\[
  \ctx \vdash e : T.
\]
The names of types and functions may suggest intended behavior, but the type
alone is the specification for this task. For example, an empty Church list
is a valid result for a type ending in \code{List a}. An implementation of
\code{sort} that always returns an empty list satisfies that type, although it
does not implement sorting. Behavioral specifications, preservation of list
length, or equivalence to an existing function require additional constraints
and separate tests.

The implementation retains three distinct notions of evidence:
\begin{enumerate}
  \item \emph{Search evidence}: a backend produces a candidate through its
    checked representation and records any remaining constraints.
  \item \emph{Target-language evidence}: GHC or the Lean elaborator and kernel
    accept the emitted implementation at the requested target type.
  \item \emph{Negative evidence}: the historical complete propositional
    fragment of Djinn may establish non-inhabitation. A failed bounded
    higher-rank search does not establish it.
\end{enumerate}
The Church harness requires the first two for every selected case. It never
treats a timeout, an empty candidate list, or a search budget as a proof that
the signature is uninhabited.

\subsection{What practical completeness means here}

The concrete acceptance obligation covers every signature in the Church
source, including the local annotations. It also motivates general rules for
arbitrary nesting, correlated instantiation, polymorphic arguments, and
capture-safe substitution. The implementation does not identify cases by
source function name or return a table of implementations.

No implementation can promise both termination and complete answers for all
System~F inhabitation requests: that decision problem is undecidable
\cite{inhabitation}. Curry-style System~F type reconstruction and type checking
are also undecidable \cite{wells}. This does not prevent independently checking
an explicitly elaborated candidate, and it does not prevent useful bounded
search. It does require precise reporting of the boundary between a checked
success and an inconclusive search.

There are also ordinary language boundaries. GHC's
\code{ImpredicativeTypes} extension enables Quick Look inference and permits
quantification beneath type constructors and explicit polymorphic type
applications. It does not make type-class instance arguments or implicit
parameters impredicative \cite{ghc}. Djex preserves the target language's
restrictions instead of treating the extension as a license to manufacture
arbitrary dictionaries.

\section{Rank-N types and impredicative instantiation}

\subsection{Quantification inside a function type}

Ordinary prenex polymorphism places type quantifiers at the outside:
\begin{lstlisting}
identity :: forall a. a -> a
\end{lstlisting}
Higher-rank polymorphism permits quantifiers inside the type:
\begin{lstlisting}
consumeIdentity :: ((forall a. a -> a) -> result) -> result
\end{lstlisting}
An inhabitant is
\begin{lstlisting}
consumeIdentity consume = consume (\x -> x)
\end{lstlisting}
The callback expects one value that can be used at every type. Replacing its
bound variable with an ordinary unification variable would weaken this
requirement: a single use at one monotype would then be mistaken for a
polymorphic value.

Introduction and elimination therefore require different operations. To
construct a polymorphic value, the search opens its quantified variables as
fresh rigid identities and constructs its body uniformly. To use an available
polymorphic value, the search chooses an instantiation of its quantified
variables. Each use may choose a different instantiation.

\subsection{Selecting a polymorphic type as an argument}

Predicative instantiation chooses a monotype for a quantified variable.
Impredicative instantiation may choose a type containing a quantifier. With
\code{Id = forall a. a -> a}, the application of a provider of type
\code{forall t. t -> t} to an \code{Id} value selects $t := \code{Id}$.
The selected type is a complete polymorphic type, not the body of a
quantifier with its scope erased.

This distinction is visible beneath ordinary type constructors:
\begin{lstlisting}
keepIdentities :: [forall a. a -> a] -> [forall b. b -> b]
keepIdentities xs = xs
\end{lstlisting}
The element schemes are alpha-equivalent. Their binder spellings differ, but
each list element remains universally quantified. Comparing textual strings
would be too strict; deleting the quantifiers would be unsound.

Quick Look is an influential example of using local structural information to
choose impredicative instances without replacing an entire production type
inference engine \cite{quicklook}. Djex uses a related engineering principle:
the shapes already present in a query are valuable evidence for selecting
instances. Its two search algorithms are not implementations of the Quick
Look algorithm, and GHC remains a separate final checker for Haskell output.

\section{Church encodings as a realistic corpus}

The basic aliases are:
\begin{lstlisting}
type Bool       = forall r. r -> r -> r
type Pair a b   = forall r. (a -> b -> r) -> r
type List a     = forall r. (a -> r -> r) -> r -> r
type Maybe a    = forall r. r -> (a -> r) -> r
type Either a b = forall r. (a -> r) -> (b -> r) -> r
\end{lstlisting}
The corpus adds \code{Tuple}, \code{Dict}, \code{Equal}, \code{LE}, and
\code{Matrix}:
\begin{lstlisting}
type Tuple a  = List a
type Dict k v = List (Pair k v)
type Equal a  = a -> a -> Bool
type LE a     = a -> a -> Bool
type Matrix a = List (List a)
\end{lstlisting}
The aliases hide substantial nesting. A dictionary contains universally
quantified pairs as elements of another universally quantified encoding. A
matrix contains polymorphic list values inside polymorphic lists. Expansion
must preserve every binder and its dependencies.

The four local annotations are the \code{interleave} helper within
\code{permutations}, \code{go} within \code{introsort},
\code{introselect} within \code{nthElement'}, and \code{go} within
\code{setOp}. The manifest records the enclosing definition and source line
for each. They are tested as generalized standalone signatures, without
captured reference implementations.

\subsection{Elimination requires a useful result type}

Consider the Church pair projection:
\begin{lstlisting}
first :: Pair a b -> a
first p = p (\x _ -> x)
\end{lstlisting}
The query requires choosing the pair's result type $r := a$. The selected
continuation then has type $a \to b \to a$. The type is not inhabited merely
by forwarding the opaque pair value; search must eliminate its quantifier.

A more revealing example is Church either elimination:
\begin{lstlisting}
eliminateEither :: (a -> c) -> (b -> c) -> Either a b -> c
eliminateEither leftCase rightCase value = value leftCase rightCase
\end{lstlisting}
Here the useful selection is $r := c$. It correlates both continuation
argument types with the desired result. Matching a provider's result suffix
against the required result supplies this information directly.

\subsection{Introduction under an argument boundary}

Currying a Church pair requires constructing a polymorphic argument:
\begin{lstlisting}
curryPair :: (Pair a b -> c) -> a -> b -> c
curryPair consume x y = consume (\continuation -> continuation x y)
\end{lstlisting}
The supplied lambda must work at every result type chosen by
\code{consume}. A search that supports only a top-level universal quantifier
cannot discharge this nested introduction obligation.

\subsection{Transporting and transforming polymorphic components}

Both engines synthesize pair exchange and pair mapping. The following forms
are alpha-renamed presentations of generated Exference candidates checked by
the whole-corpus harness:
\begin{lstlisting}
swapPair :: Pair a b -> Pair b a
swapPair p continuation = p (\x y -> continuation y x)

mapPair :: (a -> c) -> (b -> d) -> Pair a b -> Pair c d
mapPair f g p continuation =
  p (\x y -> continuation (f x) (g y))
\end{lstlisting}
The input and output result quantifiers are distinct lexical scopes. Their
relationship is established by the selected instantiation at a use site, not
by globally identifying similarly named binders.

\section{The shared representation and its invariants}

The parser-independent foundation represents variables, constructors, type
applications, arrows, products, and explicit universal quantifiers. A
quantifier retains its binder list, class context, and body. The representation
supports alpha-stable keys, free-variable analysis, substitution, and
synonym expansion. Both backends use this foundation when crossing the public
Djex API.

There are three relevant identities in the unification and instantiation
logic:
\begin{description}
  \item[Flexible variables] may be solved by the current search equation,
    subject to ownership and scope checks.
  \item[Ambient rigid variables] represent fixed types in the current
    context, including types opened by a legitimate introduction step.
  \item[Fresh comparison binders] exist only while comparing corresponding
    quantified types. They must never escape into a substitution for an older
    flexible variable.
\end{description}
A source spelling is a presentation hint; it is not sufficient evidence that
two variables have the same identity.

\subsection{Capture-avoiding substitution}

Substitution must traverse both a quantifier's context and its body while
respecting shadowing. If a replacement contains a name that would be captured
by a nested binder, the nested binder is renamed first. In particular,
\[
  [a := b](\forall b.\ a \to b)
  \quad\text{must denote}\quad
  \forall c.\ b \to c,
\]
with fresh $c$, rather than $\forall b.\ b \to b$.

The Church extractor independently applies this rule to aliases. It gives
each quantifier a fresh ASCII binder before substituting its type parameters.
The resulting expanded type is then checked against the original aliases by
GHC, so an error in this test-side transformation cannot silently redefine the
acceptance type.

\subsection{Whole polytypes remain whole}

Quantified types remain intact when they become substitution images. A
metavariable may receive a complete type such as $\forall a.\ a \to a$.
Free-variable traversal still reaches the free variables inside this type;
``opaque'' does not mean invisible to occurrence, substitution, or scope
checks. Only the bound variables are protected by their lexical scope.

\section{Djinn: structural demands become checked instances}

Djinn's propositional core uses LJT proof search. The higher-rank extension
retains quantified types as alpha-stable atoms when they must be transported
opaquely, and opens eligible positive occurrences for introduction. Checked
instantiation bridges then expose useful hypothesis instances to the
propositional search.

\subsection{One coherent view for transport and introduction}

A query may need to forward several existing polymorphic values unchanged
while constructing other polymorphic values in the same result. Treating
every positive quantifier as opaque prevents the constructions; opening every
one loses the exact types needed for forwarding. Enumerating subsets of
positive sites introduces an artificial boundary at the next omitted
combination.

The transport-directed view makes the decision at each site in one traversal.
A positive quantified scheme remains opaque exactly when an alpha-equivalent
source scheme occurs at a negative position in the query or an actual
provider. Otherwise the positive occurrence is introduced with fresh rigid
variables. For example:
\begin{lstlisting}
type Seven r = forall a b c d e f g.
  (a,b,c,d,e,f,g) -> r
type Id = forall x. x -> x

mix :: Seven r1 -> Seven r2 -> (Seven r1, Seven r2, Id, Id)
mix p q = (p, q, \x -> x, \x -> x)
\end{lstlisting}
The supplied \code{Seven} values retain their complete schemes, while each
identity result opens for construction. The same decision scales to sixteen
or twenty-four mixed sites without enumerating the corresponding subset
frontier.

Loaded consumers must use the same view and namespace as the query. The
environment therefore compiles query-aware argument views for real named
providers; these views do not introduce a new provider or assume a missing
implementation. Core search builds one coherent structural sequent, plus a
nominal version when relevant, and checks every resulting proof in the normal
proof environment.

For queries with at most eleven positive sites, this plan runs first only
when its goal formula or an active provider's \emph{symbol--formula pair} is
absent from the established cached frontiers. A target-only premise is not
evidence of such novelty. For larger queries, exact novelty enumeration itself
can be costly, so the coherent view may run early without that exhaustive
comparison. Eleven is an acceleration threshold, not a rank or site-count
limit. With no available quantified scheme the extra plan is omitted; when
the bounded novelty comparison proves the view already cached, the established
alternative streams remain in place. The richer directed-instantiation family
still waits until the earlier families and transport view have produced no
inhabitant or definitive negative evidence. These decisions affect scheduling
and useful representations, not source arity or the meaning of proof checking.

\subsection{Constructing an argument at an impredicative instance}

Selecting a polytype is only the first half of many practical examples. A
provider may require a fresh value of that polytype, rather than accept one
already present in the context:
\begin{lstlisting}
buildComposition ::
  (forall a. a -> f a) ->
  (forall a. g a -> h a) ->
  (forall a. h a -> k a) ->
  f (forall a. g a -> k a)
buildComposition pack gh hk = pack (\x -> hk (gh x))

buildNested :: (forall a. a -> f a) ->
  f (forall b. b -> f (forall c. b -> c -> b))
buildNested pack = pack (\x -> pack (\_ _ -> x))
\end{lstlisting}
In the first example, a new rigid type parameter belongs to the argument
being constructed. Both transformation providers must be instantiated at
that same parameter. In the second, the inner construction depends on the
outer argument's parameter and value. Flattening their scopes together would
admit escaping variables; isolating every inner scope from its parent would
discard this valid implementation.

Djinn's construction family therefore translates instantiated provider bodies
at the polarity of an available premise. Quantified argument obligations open
with private rigid identities. Each root construction owns an isolated family
that can specialize the original checked providers within that scope and
discover further obligations exposed by those instances. Descendants retain
their actual scope dependencies; their private variables do not become
general candidates in the enclosing query. One family shares the operational
allowance of 512 expansion attempts after selecting its root and 64 retained
premises including that root. Initial root proposals use the ordinary bounded
instance-family builder separately.

Exact schemes retained for real loaded providers are included alongside
quantified hypotheses discovered in the query, including schemes exposed
only after applying preceding ordinary arguments. For a polymorphic goal
requiring such construction, a small context runs after the primary search
opportunity and before broad historical loaded-instance guesses. It
contains the actual named scheme premises and one complete checked
construction family. This ordering prevents unsuccessful older guesses from
using the entire choice budget before a useful construction can be tried.
For newly eligible sources exposed after ordinary arguments, an already
available exact opaque source for the whole goal keeps the established
forwarding route ahead of speculative construction. Leading schemes retain
their existing policy. For a monomorphic goal, the deferred phase also tries
exact forwarding first.
No scheduling choice adds an unchecked source
assumption or changes the authority of negative evidence.

Scope ancestry follows the free variables of the original scheme and its
selected type arguments. An independent later quantified result can begin a
new scope; a dependent nested construction retains its deepest owner and that
owner's ancestors. The reuse check rejects returning to an ancestor-active
construction with its fixed private identities. Independent sibling uses of
the same provider remain available: constructing a pair of polymorphic
identities does not require an artificial sharing restriction.

The ordinary proof checker runs before the additional scope and reuse checks,
and rejected alternatives do not refund search fuel. Kind checking temporarily
renames exactly the owned private identities to fresh legal source names,
while retaining the original identities in formulas and evidence. A private
name outside the explicit ownership set is rejected. These steps provide
usable construction scopes without treating a skolem as an unrestricted
inhabitance assumption.

\subsection{The shape of an instantiation bridge}

For an available scheme
\[
 S = \forall a_1\ldots a_n.\ B,
\]
and a checked selection $\theta = [a_i := U_i]_{i=1}^{n}$, an elimination
bridge relates the scheme to its specialized body $B\theta$. The bridge
retains the original source scheme and the complete argument vector.
Translation, kind checking, proof checking, and generated-term association
remain responsible for admitting the instance.

The new directed family is additive. It does not change the meaning of
existing opaque transport or polarized introduction rules. It also does not
turn an unproductive bounded rule into a negative proof.

\subsection{Match useful shapes instead of guessing all tuples}

Let $\mathcal D$ be the deduplicated collection of eligible demands from the
elaborated query and its retained type atoms. A demand may contain ambient
variables, arrows, constructors, or a complete quantified subtree. Every free
variable in it must belong to the allowed current scope.

The directed matcher compares the entire body $B$ and each function-result
suffix of $B$ with each demand in $\mathcal D$. Only the leading variables
$a_1,\ldots,a_n$ are selectable in the resulting leading-argument vector.
Ambient source and demand variables remain rigid. When a result spine exposes
a later context-free universal group, its temporary binders may participate
in matching the eventual result; their solutions are discarded from the
current vector. Ordinary checked application and elimination must still reach
and instantiate that later group. Repeated occurrences of a selectable variable must
produce alpha-equivalent selections. Source and demand binders are first
made lexically unique using legal reserved source names, so identically
spelled nested binders cannot be mistaken for outer parameters and ordinary
kind validation can still inspect the selected types.

For the either eliminator, the scheme body is
\[
 (a\to r)\to(b\to r)\to r.
\]
The final result suffix $r$ matches the demand $c$, producing the correlated
selection $r:=c$. This avoids spending the choice-point budget on irrelevant
instances before reaching the useful continuation type.

For a wide provider, one traversal can solve twelve or more correlated
variables. The directed family's eligibility has no fixed six-binder cutoff.
That statement is specific to this family: historical tuple-enumeration and
provider-evidence routes retain their separately documented limits.

\subsection{Nested quantifiers and scope protection}

When a pattern and demand both contain quantifiers, the matcher establishes a
lexical correspondence between their binders and compares the associated
contexts and bodies. The correspondence takes precedence over outer
selectable names. A selected image may not contain a demand-side binder
introduced only by this nested comparison.

Thus a pattern $\forall b.\ a\to b$ may match
$\forall c.\ \code{Int}\to c$ with $a:=\code{Int}$. It must not match
$\forall c.\ c\to c$ by selecting $a:=c$, because $c$ is not available outside
that comparison. Alpha-renaming and scope checking are separate duties.

\subsection{Unobserved choices and operational bounds}

A result shape need not mention every quantified parameter. The remaining
parameters draw from the same checked vocabulary. Repeated, or diagonal,
selections are visited before the full Cartesian tail, which prevents an
unproductive first coordinate from starving a useful repeated choice.
Fully determined matches are emitted first. Partially determined matches
then take turns contributing vectors, so guessing parameters for one match
cannot indefinitely suppress another match's useful first selection.

Search-plan scheduling also matters. After the structural and nominal plans,
focused plans try individual checked instantiation axioms to accelerate the
first inhabitant. They are skipped once an earlier plan has produced a
candidate. Rich compound demand-derived axioms are kept separate from the
historical correlated family and act as an inhabitation fallback only when
the earlier search has no result. This preserves established alternative
ordering and cutoff behavior for already inhabited goals and avoids opening
a large new alternative stream after a satisfactory result has been found.
This policy prioritizes the first checked inhabitant under the configured
limits; it does not guarantee success for every inhabited type or enumerate
every typable program.

The current family builder retains at most 16 axioms per scheme and 64 per
family, with a 512-attempt bound. These are operational limits, not language
rank limits. The directed rule remains positive-only. A search that misses an
inhabitant because of these limits must remain inconclusive.

The raw proposed-vector allowance is applied before alpha-equivalence
deduplication. Otherwise a request could traverse an arbitrarily large sequence of
duplicate or rejected proposals while apparently retaining only a few
instances. A scheme whose allowance is exhausted stops before forcing the
remaining lazy proposal stream. This bounds the number of raw vectors inspected
before deduplication and the number of retained instances. It does not impose
a wall-clock or LJT choice-point bound; those are separate query limits, and
an unbudgeted provider search may remain expensive.

The main implementation locations are
\path{djinn/src-core/Djinn/Internal/DirectedInstantiation.hs},
\path{djinn/src-core/Djinn/Internal/Instantiation.hs}, and the checked query
integration in \path{djinn/src-core/Djinn/Core.hs}.

\subsection{Exact argument vectors supplied by a frontend}

A frontend may already know the complete correlated type arguments of a
retained provider. Such evidence is different from a heuristic request to
guess arguments from a scalar vocabulary. Both backend adapters accept exact
vectors through the shared provider-instantiation APIs, in forms with inferred
argument kinds or explicitly supplied ground kinds.

The vector's required length is now derived from the validated retained
provider's leading binder count $n$. Validation observes at most $n+1$ cells
of the caller's argument-list spine before traversing any argument contents.
A short vector fails; an extra cell proves that a long or cyclic vector has
the wrong arity. Thus a twelve-binder provider can receive twelve arguments
without making infinite caller lists a termination hazard.

The exact-evidence route still limits a request to 32 assignment vectors, and
each explicitly supplied kind has a 129-node structural preflight bound.
Provider identity, exact argument count, substitution consistency, kinds,
complete specialized type, and candidate association all remain checked.
The historical six-argument constant continues to bound heuristic Cartesian
enumeration; it is no longer a language bound on exact supplied vectors.

The facade regression suite exercises eight- and twelve-argument vectors for
both engines, through both inferred-kind and explicitly kinded APIs, including
GHC replay of the emitted applications. It also checks oversized and cyclic
vectors. These tests distinguish a wider legitimate provider from malformed
external evidence, instead of simply increasing a global magic number.

Exact frontend assignments have dedicated priority plans as well as combined
instance-family contexts. The historical families retain their relative
order, while the coherent transport and loaded-construction accelerators may
interpose when no earlier candidate has succeeded. Priority is an execution
policy; every route retains the same evidence and target-language checking
requirements.

\section{Exference: structural equations between quantified types}

Exference already introduces nested universal goals with branch-local rigid
variables and instantiates scoped providers at independent use sites. The
new unifier extends equations between quantified types while retaining whole
polytypes as possible substitution images.

\subsection{A motivating equation}

Suppose a candidate must reconcile
\[
 \forall a.\ a\to \beta
 \quad\text{and}\quad
 \forall b.\ b\to \code{Int}.
\]
Pure alpha-equivalence is insufficient: the free flexible variable $\beta$
still needs to be solved. The correct solution is $\beta:=\code{Int}$.
Conversely,
\[
 \forall a.\ a\to\beta
 \quad\text{and}\quad
 \forall b.\ b\to b
\]
must not solve the older $\beta$ with the newly opened $b$.

\subsection{Open corresponding binders in a private namespace}

The unifier maintains a tagged representation distinguishing flexible
variables, ambient rigid constants, and private bound-comparison identities.
For an equation between two quantified schemes, it opens one corresponding
binder from each side with the same fresh private rigid identity. Repeating
this operation handles both grouped and successively written quantifier
prefixes.

The remaining contexts and bodies become structural equations. Contexts are
compared for the same class, arity, and argument structure. This operation
does not solve a class constraint or create a dictionary; those obligations
belong to the existing scoped constraint solver.
Grouping and successive opening are supported for context-free prefixes and
corresponding context layers. Moving a constraint across a quantifier boundary
or treating differently ordered constraints as interchangeable is not part
of this structural equation rule.

\subsection{Occurrence and escape checks traverse polytypes}

Before a variable is bound, the unifier checks:
\begin{enumerate}
  \item the variable does not occur free in its own proposed replacement;
  \item no private comparison binder occurs in that replacement;
  \item substitution into outstanding equations and existing bindings remains
    capture-avoiding and structurally valid.
\end{enumerate}
The same checks apply when a variable occurs inside a quantified substitution
image. Updating all equations and already recorded substitutions prevents an
indirect chain of bindings from hiding an escape.

The first condition rejects a cyclic solution such as
$\beta := \forall a.\ a\to\beta$. The second rejects
$\beta:=\forall c.\ b\to c$ when $b$ exists only inside the current
quantified comparison. A complete closed image such as
$\beta := \forall c.\ c\to c$ remains admissible.

The implementation is in
\path{exference/src-core/Language/Haskell/Exference/Core/Unify.hs}.
Candidate reconstruction and expression checking remain independent consumers
of the selected substitutions. A unification success is not itself a
target-language acceptance result.

\subsection{Forwarding a whole scheme to a flexible argument goal}

Structural unification is useful only when search presents the right equation
to it. A provider may expect an argument whose type is still flexible. If a
polymorphic value is available, eagerly eliminating all of its quantifiers
loses the opportunity to choose the entire scheme as that argument's type.

Exference therefore retains a sibling search branch that forwards the
available complete polytype to the flexible argument goal. The ordinary
instantiated-use branch remains available. A representative shape is:
\begin{lstlisting}
apply :: forall a r. a -> (a -> r) -> r
identity :: forall b. b -> b
consumeIdentity :: (forall b. b -> b) -> result

-- The argument metavariable of apply can select the complete Id scheme.
use :: result
use = apply identity consumeIdentity
\end{lstlisting}
The provider declarations in this example describe a contextual synthesis
environment. The search result need not instantiate \code{identity} to a
monotype before selecting it for \code{apply}. Its complete scheme supplies
the correlated choice, and the ordinary expression checker validates the
result. This change is in the Exference search module
\path{Exference/Core/Internal/Exference.hs}; it does not weaken the checker.

\subsection{A polymorphic provider exposed by a term application}

Quantifiers can occur after ordinary value arguments, so useful provider
schemes need not be visible at the beginning of an expression. Consider:
\begin{lstlisting}
interleaved :: forall seed.
  seed -> (seed -> (forall a. a -> F a)) -> F (forall b. b -> b)
interleaved seed maker =
  let localFactory = maker seed
  in localFactory (\x -> x)
\end{lstlisting}
The first application exposes the scheme
\code{forall a. a -> F a}. Its next use must select the identity polytype
and construct a polymorphic identity argument. The seed may have any type;
neither a distinguished unit value nor a special factory name is required.
GHC independently accepts this implementation at the displayed signature.

Exference's general partial-application rule already creates a local binding
for the residual type and continues search in the extended lexical scope.
The independent checker must preserve that same information boundary. A
bottom-up inference pass over the let body could infer the identity lambda
as a monotype before discovering that the expected result is
\code{F (forall b. b -> b)}. Let checking therefore checks the binding at
its retained annotation and checks the body at the trusted outer expected
type. The annotation is independently validated; it is not an unchecked
instruction to accept the desired result. Regressions exercise one and two
ordinary arguments, nominal and ambient seed types, and rejection of an
unrelated rigid annotation.

\subsection{Successive term and quantifier layers}

The same issue can repeat along one provider spine:
\begin{lstlisting}
type Id = forall c. c -> c

alternating :: forall seed.
  seed ->
  (seed -> forall a. a -> seed -> forall b. b -> G a b) ->
  G Id Id
alternating seed maker =
  let secondFactory = maker seed @Id (\x -> x) seed
  in secondFactory (\x -> x)
\end{lstlisting}
The current instance of \code{a} must be selected before a later result
quantifier exposes the final \code{G a b}. Exference uses a positive
lookahead branch: it temporarily opens successive residual forall layers,
peels their value arrows, and matches the ultimate result against the current
demand. Only substitutions for the provider's current leading binders are
retained. A retained image containing a fresh variable belonging solely to a
future layer is rejected. Every value argument and every later provider use
must still be synthesized and independently checked in its actual scope.
Lookahead is an instance-selection heuristic, never a proof of the goal.

Source reconstruction must also preserve the instantiation boundary. The
fully inlined term \code{maker seed (\textbackslash x -> x) seed
(\textbackslash y -> y)} is accepted by GHC at this signature. In contrast,
introducing two inferred let bindings without a type application can infer
the first identity as a monotype and reject the final result. Applying
\code{@Id} to an inferred let alias is also invalid, because that inferred
binder is not a specified type-application slot.

The generated form above keeps the selected type application on the original
applied-expression spine. Nonrecursive aliases used under explicit type
applications are selectively inlined with the shared capture-avoiding
substitution; other lets remain useful expected-type boundaries. The
normalized expression is checked independently, as is any later repaired
expression after the same normalization. No extra source declaration or
unchecked type assertion is introduced. Ambient images use the scoped type
patterns described below and receive the same final GHC check. Selective
inlining can change sharing and runtime cost; type correctness does not certify
an optimal evaluation strategy.

\section{Haskell output and independent compiler checking}

The Haskell acceptance pipeline is:
\[
 \text{type and context}
 \longrightarrow \text{checked request}
 \longrightarrow \text{candidate}
 \longrightarrow \text{rendered source}
 \longrightarrow \text{GHC acceptance}.
\]
The library independently checks its internal candidate evidence; an ordinary
public API call does not itself launch GHC. The corpus harness or the library's
consumer supplies the external compiler replay shown in the final step.
Each arrow has a distinct failure mode. A representable internal type may
require an explicit annotation in source. A correct quantified instance can
be rendered incorrectly if its binder scope is lost. GHC's inference of an
unannotated expression may choose a different instance from the one intended
by search. The final compiler check detects these failures.

\subsection{Nested instantiation can require explicit type evidence}

Quick Look can infer many impredicative applications from their argument and
result shapes, but that inference is not an unrestricted structural unifier
beneath arbitrary universal quantifiers. Consider an abstract type constructor
\code{F :: Type -> Type} and the query:
\begin{lstlisting}
nested :: (forall a. F (forall b. b -> a))
       -> F (forall c. c -> (forall d. d -> d))
\end{lstlisting}
The correct choice is $a := \forall d.\ d\to d$. The source
\code{nested p = p} nevertheless fails GHC's inference at this type. Explicit
evidence records the successful choice:
\begin{lstlisting}
nested p = p @(forall d. d -> d)
\end{lstlisting}
This example separates successful internal instantiation from successful
source reconstruction. Returning the same unannotated variable after proving
that it can be instantiated is insufficient.

The chosen polytype can also mention an ambient type variable:
\begin{lstlisting}
nestedAmbient :: forall x.
  (forall a. F (forall b. b -> a))
  -> F (forall c. c -> (forall d. x -> d -> d))

nestedAmbient p = p @(forall d. _ -> d -> d)
\end{lstlisting}
The anonymous placeholder preserves the quantified shape while allowing GHC
to recover the ambient \code{x} from the expected result. It avoids assuming
that an internal fresh variable has the same textual name as the caller's
scoped type variable. GHC 9.12.4 accepts these interior placeholders in visible
type applications without requiring \code{PartialTypeSignatures}. They are
inference variables checked against the original signature, not unchecked
type casts or values supplied to synthesis.

The distinction between locally bound and ambient variables remains crucial.
The explicit \code{d} above belongs to the polytype and must retain its binding
relationship. Repeated ambient variables, two correlated ambient variables,
and higher-kinded ambient variables such as \code{f} in
\code{forall d. f d -> d -> f d} exercise this obligation further. An
application such as \code{@(forall d. \_ d -> d -> \_ d)} still requires GHC to
reconcile every placeholder with the same requested result type. The target
compiler is the final authority on whether that reconstruction succeeds.

The shared generated-expression representation separates a fully inferred
argument, a specified closed type, and a partially specified type pattern.
The pattern preserves alpha-normalized lexical identities for all inner
binders and uses anonymous holes only for ambient variable occurrences. The
exact correlated instantiation remains part of the checked evidence; the
pattern is source syntax derived from that evidence. The Exference expression
checker reconstructs fresh metavariables for the holes and checks the whole
application against its actual expected type. A pattern containing holes is
not by itself a certificate that any instantiation is legal.

The required choice may become known only when an argument is supplied:
\begin{lstlisting}
delayed :: (forall a. F (forall b. b -> a) -> result)
        -> F (forall c. c -> (forall d. d -> d))
        -> result
delayed p q = p @(forall d. d -> d) q
\end{lstlisting}
Here the result type alone does not determine \code{a}; checking \code{q}
does. The unannotated source \code{p q} fails the same GHC inference boundary.
Reconstruction therefore has to inspect the fully solved use, including its
arguments, rather than only the substitution discovered while matching the
provider's result. This obligation applies equally to scoped function
arguments and named global providers, including when the selected polytype
contains ambient variables.

Exference handles that timing by recording every ordinary local and global
occurrence during the successful checker traversal, independently of optional
term-graph construction. Abandoned checker alternatives roll back their
traces. After solving the complete expression, the final substitutions are
applied to the retained occurrence types and the trace is aligned with the
original expression. A source-order mismatch rejects the candidate. Existing
explicit global application spines retain their evidence instead of being
decorated twice.

Only the necessary visible prefix is inserted: it ends at a selected binder
whose polymorphic image cannot be inferred through its exclusively nested
quantified occurrences. Earlier positions can use \code{@\_} to preserve the
source binder order. The entire rewritten expression is independently checked
again before publication. This makes the completed-use repair an elaboration
step supported by fresh checking, rather than a textual substitution applied
to already printed code.

\subsection{The source inventory is independent of synthesis}

\path{test-church/extract_corpus.py} enumerates source signatures and aliases,
then invokes GHC with:
\begin{lstlisting}
ghc -v0 -fno-code -fno-write-interface -ddump-types \
    -fprint-explicit-foralls -dppr-cols=10000 \
    docs/examples/Church.hs
\end{lstlisting}
The original source typechecks under GHC 9.12.4. Its bodies are used only to
resolve partial signature holes. For example, GHC resolves the omitted
function type in \code{uncurry}, and the resulting full signature is the
query. The extractor does not export a source expression to either backend.

The manifest records the source SHA-256, compiler version, source line,
original signature, full resolved type, expanded AST, classification, and
accepted assumptions. Its companion TSV is a compact projection consumed by
the Haskell runner. The runner first regenerates the expected inventory in
memory and rejects stale checked-in artifacts.
The source hash is computed from UTF-8 text without a BOM and with line
endings normalized to LF, so Windows and Unix Git checkouts have the same
corpus identity.

\subsection{Checking both sides of alias expansion}

For each candidate, the generated compiler module contains two declarations:
\begin{lstlisting}
candidate :: FullyExpandedQuery
candidate = generatedExpression

candidate_original :: OriginalAliasPreservingType
candidate_original = candidate
\end{lstlisting}
The first checks the exact expanded query. The second asks GHC to reconcile
that query with the original aliases. For partial cases both declarations
contain the same explicit default argument. The generated module enables
\code{RankNTypes}, \code{ImpredicativeTypes},
\code{ScopedTypeVariables}, and \code{TypeApplications}.

For each inherently partial Haskell case the compiler fixture also checks:
\begin{lstlisting}
candidate_partial :: ExactOriginalSignatureWithoutDefault
candidate_partial = candidate undefined
\end{lstlisting}
The synthesized helper's first argument is exactly the declared element
default. The wrapper supplies bottom at that one boundary after synthesis;
it does not expose a universal bottom provider to any search.

The module uses \code{NoImplicitPrelude}. Its ten Church aliases are copied as
type declarations only. The search environment has an abstract \code{Int}
type, and the 16 classified concrete-provider queries additionally have
\code{churchZero :: Int}. No other named implementation, bottom value, or
recursion helper is a candidate provider. The compiler module imports
\code{undefined} solely for the explicit partial wrappers.

\section{Partial signatures and inherently partial functions}

Two different uses of the word ``partial'' must be kept separate. A
\emph{partial type signature} contains a source hole such as \code{\_}; GHC
can resolve that hole from the program. An \emph{inherently partial function}
has an original type with no total closed inhabitant in the intended
parametric setting. Resolving a signature hole does not supply a missing
element value.

\Needspace{9\baselineskip}
\begin{proposition}
There is no total closed implementation of
\code{forall a. List a -> a} under the pure Church interpretation.
\end{proposition}
\begin{proof}
Choose an empty element type $E$. The Church empty list
$\lambda c.\lambda z.z$ inhabits $\code{List}\ E$. Applying the purported
total implementation at $E$ to this list would produce an element of $E$,
which is impossible.
\end{proof}

The same argument applies to \code{fromJust} by supplying Church nothing, to
\code{fromLeft} and \code{fromRight} by selecting the other alternative, and
to dictionary lookup by supplying an empty dictionary. Extremum and
seedless-fold signatures still admit empty input containers; the signatures
do not express their informal nonempty precondition.

The explicit policy is to expose a default of exactly the required element
type:
\begin{lstlisting}
headWithDefault :: forall a. a -> List a -> a
fromJustWithDefault :: forall a. a -> Maybe a -> a
minimumWithDefault :: forall a. a -> LE a -> List a -> a
\end{lstlisting}
This is a contextual total implementation. It makes no claim that the original
default-free signature has become total. Supplying an unrestricted
\code{forall a. a} value would make every type trivially inhabitable and
would invalidate the intended acceptance test; the harness does not do so.
Haskell's exact original type is nevertheless checked by wrapping the
synthesized helper with an explicit partial value at that single default
argument. Lean keeps the stated default premise and introduces no bottom.

The 19 cases are \code{head}, \code{last}, \code{at}, \code{foldl1},
\code{foldr1}, \code{fromJust}, \code{fromLeft}, \code{fromRight},
\code{maximumBy}, \code{maximumOn}, \code{minimumBy}, \code{minimumOn},
\code{minMaxBy}, \code{minmaxElement}, \code{atKey}, \code{reduce},
\code{maximum}, \code{minimum}, and \code{minMax}. The manifest retains the
original signature alongside the augmented query for every one.

\section{Lean integration and its universe boundary}

Leant reifies supported Lean types to the shared synthesis fragment, invokes
the Haskell search backends, renders candidate Lean terms, and checks them
against the original Lean goal. The result must pass Lean elaboration and
kernel checking. Reification and Haskell search are not trusted replacements
for this final check.

\subsection{Lean Type is predicative}

Haskell's impredicative type universe cannot be copied literally into Lean's
\code{Type} hierarchy. Lean has \code{Type u : Type (u + 1)}, its universes are
not cumulative, and universe-polymorphic declarations are instantiated at
specific levels. \code{Prop} has its own impredicative behavior; it is not a
substitute for arbitrary computational data \cite{leanuniverses}.

For example, if $a : \code{Type}\ u$, a Church list that eliminates into
$r : \code{Type}\ v$ has the form
\[
 \Pi(r:\code{Type}\ v).\ (a\to r\to r)\to r\to r.
\]
Its universe is \code{Type (max u (v + 1))}. Nesting such a list as an element
of another encoding therefore imposes real universe constraints. It is
incorrect to force every quantified type variable to \code{Type 0} and then
assume every Haskell impredicative instance remains valid.

\subsection{Universe-correct counterparts of the corpus}

The Lean corpus generator consumes the expanded AST and translates every
quantified type variable to an explicit binder of type \code{Type \_}. Arrows
and applications retain their structure, and Haskell \code{Int} maps to Lean
\code{Int}. Each universe metavariable is constrained by elaboration of the
goal and generated term.

This checks a valid universe instantiation of each Haskell signature. It does
not claim that one implementation works for all independent assignments of
universe levels, nor that Lean's computational universes are impredicative.
The 19 default-bearing cases insert exactly the same element-default premise
as the Haskell harness.

The frontend has explicit wire-format resource limits. Quantifier metadata
and class-context argument lists each have their own 128-entry bound, while
individual kind shapes retain the 129-node bound. Class argument width is
independent of the old six-choice heuristic. The producer derives an exact
provider vector's arity from its source binders, under its existing serializer
depth fuel of 80, and the receiving backend validates that vector against the retained
provider. These transport limits must be reported as such; they are not
grounds for a mathematical non-inhabitation claim.

\Needspace{8\baselineskip}
One small counterpart is:
\begin{lstlisting}
example : (forall (a : Type _) (b : Type _),
    (forall r : Type _, (a -> b -> r) -> r) -> a) :=
  fun a b p => p a (fun x _ => x)
\end{lstlisting}
An abstract provider can also select a polymorphic value at a suitable higher
universe:
\begin{lstlisting}
example : (forall F : Type _ -> Type _,
    (forall a : Type _, F a) ->
      F (forall b : Type _, b -> b)) :=
  fun F supply => supply _
\end{lstlisting}
These examples are included in \path{Leant/test-church/UniverseProbe.lean}.

\subsection{Replay and the prohibition on invented proof evidence}

The Lean acceptance runner disables the general library provider inventory
and classical fallback for pure and default-bearing cases. Integer-provider
cases enable the explicitly defined zero and the ordinary available integer
constructors. It records the exact accepted term for every query, writes all
terms into standalone Lean declarations, and invokes Lean again on that file.
The replay also inspects reported axioms for \code{sorryAx}.

An emitted term must therefore survive both interactive verification and
standalone kernel replay. A declaration containing \code{sorry}, a forged
axiom, or a missing query is a failure. A successful replay of 349 candidates
does not compensate for a missing 350th candidate.

\subsection{Implicit type binders are lambda-group boundaries}

Lean output has a subtle additional scope obligation. Adjacent value lambdas
can often be printed in one group, but an intermediate implicit universal
quantifier can prevent that regrouping. Consider a callback that accepts a
polymorphic identity and returns another polymorphic identity:
\begin{lstlisting}
forall R : Type,
  (((forall {A : Type}, A -> A) ->
      (forall {B : Type}, B -> B)) -> R) -> R
\end{lstlisting}
The following candidate preserves the intermediate quantifier boundary:
\begin{lstlisting}
fun _ f => f (fun _ => fun x => x)
\end{lstlisting}
Flattening its inner groups to \code{fun \_ x => x} can make Lean assign the
binders at the wrong expected-type boundary. Leant's renderer reconstructs
lambda groups from the reified type fragment and starts a fresh group at each
intermediate implicit quantifier. Initial implicit-binder elision retains its
existing behavior. This is a source-rendering correction, independently
checked by the Lean examples and kernel replay, rather than a new logical
search rule.

\subsection{Expected types must reach nested arguments and let aliases}

An unannotated value lambda does not reveal whether its input is an ordinary
monotype or a value with its own universal quantifier. For example, an
application of a provider of type \code{forall a. a -> F a} to an identity
lambda can require a polymorphic argument because the expected application
result is \code{F (forall b. b -> b)}. Leant matches that expected result
against the provider result and propagates the resulting capture-aware type
substitution back into every argument domain. Lambda reconstruction then sees
the necessary explicit or implicit type binders before emitting the argument.

An explicit closed structural type argument also supplies instantiation
evidence at its original source binder before the final result expectation
becomes available. This matters when a provider is partially applied, selects
one or two polymorphic types, and then returns another quantified function.
The renderer uses \code{visibleTypeArgumentClosedType} for this evidence and
preserves the remaining quantifier layers. A partially specified argument
containing anonymous holes is still rendered as a pattern, but is not treated
as exact closed evidence.

A plain let alias retains its exact type, including universal quantifiers
exposed after the preceding term arguments. Its body is fitted against the
enclosing expected result. Silently opening that retained quantifier too early
would hide the very boundary at which a later application selects a
polymorphic type. The same distinction is relevant when an argument result
contains unresolved variables introduced by opening a universal quantifier:
such a result supplies no reliable type constraint for the caller. The caller's
expected argument type must instead guide reconstruction. Actual ambient free
variables remain rigid throughout this process.

When an applied let alias still lacks an exact fitting domain, the renderer
also considers a capture-avoiding normal form that inlines single-use let
bindings. It refits that expression from the original goal; it does not
duplicate a shared computation or perform eta contraction. The established
renderings remain the first complete prefix. Each normal form retains its
own bound of three render styles times 32 assignment choices, for at most
192 candidates across the two forms. A successfully fitted alternative can
also recover from a primary rendering failure after strict source admission.

Visible selections at an implicit quantifier reached after an ordinary
argument require positional syntax at the original known head. For example,
Lean accepts \code{@factory seed ...}; it does not accept the analogous
\code{@(factory seed) ...}. The renderer follows the complete checked source
spine from that head and exposes intervening implicit type and instance slots
with explicit arguments or inference placeholders. Exact named-provider
evidence retains its separate metadata path: if flattening the complete spine
also exposes a leading retained vector, its original forall-domain and binder
visibility metadata are still used. An inconsistent retained vector fails
instead of falling back to reconstructed evidence.

An open visible type annotation spells its known universal binders explicitly,
even when the target type writes corresponding binders implicitly. A fitting-only
visibility overlay aligns those known structural positions with the printed
annotation after expected-type resolution. It preserves inferred type
identities and the result fragment, leaves wildcard subtrees unchanged, and
does not override retained exact provider vectors.

A closed annotation can similarly use explicit forall syntax while the
expected result uses implicit binders. Final-result matching therefore ignores
only the explicit/implicit flag of corresponding universal binders. It still
checks the complete structure, rigid variables, correlations, and scope, and
preserves the original result and annotation metadata. Known-argument matching
and retained provider-evidence matching keep their stricter visibility checks.
This separates the information needed to infer an application result from the
information needed to print its value arguments correctly.

There is one further inference boundary when an argument's source type
variable is still unresolved and another universal group occurs later in the
provider's result spine. Lean may elaborate that argument before it knows
that the argument must itself be polymorphic. For an implicit polymorphic
identity, the renderer then writes the implicit type lambda explicitly:
\begin{lstlisting}
factory seed _ (fun {_} value => value) seed _ ...
\end{lstlisting}
This supplies the lambda's binder shape without guessing its type. The same
operation follows nested implicit groups inside the argument. It is restricted
to unresolved source domains before a later quantified layer; already known
annotation domains keep their existing rendering. A private marker introduced
after name uniquification prints the anonymous implicit binder and has no
term-variable occurrences. It creates no new type identity or provider and
does not add a rendering cohort. The eight constructive declarations in
\path{Leant/test-church/MixedSpineSyntax.lean} replay with empty axiom inventories.

Constructor-pattern reconstruction uses the actual scrutinee type to determine
its field types. It cannot safely infer those types from the order of source
names or unrelated quantified variables. Together, these corrections preserve
the information needed by Church continuations, nested optional values, and
association-list operations such as \code{lookupAll}. Every correction is
tested by replaying the exact rendered Lean term.

\subsection{Search fallback and bounded backend communication}

Some Church goals return a quantified continuation after ordinary inputs that
the matching implementation may ignore. If the earlier Exference search lanes
find no term, Leant also retries without multi-constructor pattern branching.
This preserves an opportunity to construct the continuation directly. The
fallback uses the same type checking and candidate verification boundaries;
it introduces no new provider or logical assumption.

Even unrelated standard datatypes can compete with the required partial
applications in a bounded queue. If all earlier lanes still miss, Leant
retries Exference with the transitive subset of the standard inventory needed
by the goal, exact type assignments, and every retained foreign declaration.
The closure follows aliases, datatype fields, class constraints and methods,
instance prerequisites and heads, and tuple identities, including unit.
All declarations supplied from Lean are preserved unchanged. The original
full variable-name table and provider ratings are reused, and the candidate's
authority records the actual smaller checked session. This is a positive
fallback after an unsuccessful full-context search; it does not add a
provider or establish negative evidence from the smaller context.

The subprocess boundary must enforce its requested time limit as well. On
Windows, an owned process job and a pipe reader have different capabilities
from their POSIX counterparts. Leant preserves ownership of the child process
job and obtains reply lines through an owned reader with a bounded queue, so a
blocked pipe read cannot silently disable the request timeout. These changes
concern execution reliability; they provide no type-theoretic evidence in
place of elaboration and kernel replay.

\subsection{Optional provider assignments have their own work budget}

Discovering a provider's ordinary type and discovering additional exact
instance-derived type arguments are separate operations. The latter is
optional: recursive instance search for a declaration such as
\code{Or.by\_cases}, whose constraints involve \code{Decidable}, must not
discard that provider or the surrounding inventory merely because useful
exact assignment evidence was not found.

Assignment discovery retains its limits of 32 inspected seed heads and sixteen
complete vectors per provider. It additionally shares one allowance of 128
instance-resolution attempts across every seed head and every recursive
constraint branch for that provider. The counter lives outside rollbackable
metavariable state, so restoring a failed branch does not restore spent work.
Each accepted vector still contains exactly the source-derived arguments and
passes the existing completeness, kind, context, and identity checks. The
attempt allowance restricts discovery work, not type arity or vector shape.

The optional computation also receives at most 2,000 additional Lean
heartbeats, clamped to the remaining outer command allowance. It never resets
or enlarges that outer budget. Ordinary Lean errors and exhaustion of the
local heartbeat allowance omit optional assignment evidence while preserving
the provider and earlier and later inventory entries. Attempt-quota exhaustion
can return the already completed vector prefix. Outer command exhaustion,
interrupts, and other runtime exceptions still propagate. Debug diagnostics in
the Haskell driver retain the actual Lean discovery error, so a failed
discovery is distinguishable from an empty successful inventory.

Five standalone Lean regressions pass for this boundary: actual discovery
through the recursive \code{Decidable} case, an ordinary error, local
heartbeat exhaustion, recursion-error propagation, and interrupt propagation.
These checks establish the behavior of the emitted Lean discovery code and
supplement the rebuilt executable's independent corpus and fixture checks
reported below.

\section{Examples from simple to demanding}

The examples below separate the structural capability exercised from any
informal behavioral expectation.

\subsection{Basic construction}
\begin{lstlisting}
pair :: a -> b -> Pair a b
pair x y continuation = continuation x y

nil :: List a
nil _ seed = seed
\end{lstlisting}
Both require introducing a result quantifier. The pair additionally captures
two outer values without specializing the result type.

\subsection{Nested continuation consumption}
\begin{lstlisting}
uncurryPair :: (a -> b -> c) -> Pair a b -> c
uncurryPair combine p = p combine

consumePolymorphic :: ((forall a. a -> a) -> r) -> r
consumePolymorphic consume = consume (\x -> x)
\end{lstlisting}
The first eliminates a supplied universal value at a chosen result type. The
second introduces a universal value inside a callback argument.

\subsection{A polymorphic payload inside an abstract constructor}
\begin{lstlisting}
-- A synthesis query; F is an abstract type constructor.
(forall a. f a) -> f (forall b. b -> b)
\end{lstlisting}
The provider must be instantiated at a complete polytype. Expanding or
erasing that polytype would change the requested result. The corresponding
Lean example above chooses consistent universe levels rather than assuming
same-level self-containment.

\subsection{Wide correlated instantiation}
\begin{lstlisting}
(forall a1 a2 a3 a4 a5 a6 a7 a8 a9 a10 a11 a12.
   (a1,a2,a3,a4,a5,a6,a7,a8,a9,a10,a11,a12))
 -> (x12,x11,x10,x9,x8,x7,x6,x5,x4,x3,x2,x1)
\end{lstlisting}
Matching the result tuple solves twelve related selections together. This is
a targeted regression for the removal of the historical fixed binder-count
restriction from the directed matching family. The premise is intentionally
very strong; this example tests provider instantiation, not the construction
of a closed inhabitant of that premise.

\subsection{Nested dictionaries and matrices}
The corpus includes signatures such as:
\begin{lstlisting}
collapse :: Equal k1 -> Equal k2 -> Dict k1 (Dict k2 v)
         -> Dict (Pair k1 k2) v

mapAccumL :: (s -> a -> Pair s b) -> s -> List a
          -> Pair s (List b)

transpose :: Matrix a -> Matrix a
\end{lstlisting}
The first has quantified encodings nested in dictionary keys and values; the
second combines a state argument, a polymorphic pair result, and a
polymorphic list result; the third moves through two layers of list encoding.
The acceptance test permits any matching inhabitant. For example, returning
an empty dictionary satisfies \code{collapse}'s type. A requirement to
actually flatten all entries must be specified and verified separately.

\section{Acceptance results and reproduction}

The source inventory has canonical UTF-8/LF SHA-256
\begin{center}\small\ttfamily
782e4edaa5bf813e30e39ae02d52278ab0566315ebc521401a947b98c44cfd11
\end{center}
and was extracted using GHC 9.12.4. The four local annotations are included in
the total-case count.

\begin{center}
\begin{tabular}{lrrr}
\toprule
Case class & Signatures & Djinn + GHC & Exference + GHC \\
\midrule
Pure total context & 315 & 315 & 315 \\
Concrete integer provider & 16 & 16 & 16 \\
Explicit element default & 19 & 19 & 19 \\
\midrule
All signatures & 350 & 350 & 350 \\
\bottomrule
\end{tabular}
\end{center}
Every Haskell success above includes both the expanded query declaration and
the alias-preserving forwarding declaration. The 19 partial cases per engine
also include an exact-original-signature partial wrapper. Thus the run checks
700 synthesized helpers through 1,438 typed declarations: 700 expanded
declarations, 700 alias-preserving forwarding checks, and 38 explicit partial
wrappers. The forwarding declarations and wrappers are checks and adapters,
not additional synthesized candidates. Every original signature is checked;
the 19 original partial signatures remain explicitly partial.

The final explicit-budget Haskell replay uses Djex implementation revision
\begin{center}\small\ttfamily
e2eb71e321c523966e47e6eedb1e1d5c14f4e7de
\end{center}
with 10,000 steps or choices and a two-second timeout per query. The runner
rechecks the source inventory before synthesis. Its generated modules,
per-query results, and compiler diagnostics are retained under
\path{test-church/results/published-layered-rankn}; a local provenance record
also identifies the tested executable by SHA-256.

From the Djex repository root:
\begin{lstlisting}
python test-church/extract_corpus.py --check
cabal test djex-church-tests --test-show-details=direct \
  --test-options="--backend both --timeout 2 --steps 10000"
\end{lstlisting}
For a development run against an already built library:
\begin{lstlisting}
cabal exec -- runghc -package=djex test-church/Main.hs \
  --backend both --timeout 2 --steps 10000 \
  --report-dir test-church/results/acceptance
\end{lstlisting}
The backslash continuations are typographical shell notation; on PowerShell,
enter the command on one line or use its native continuation syntax. The
runner defaults to both engines and records every query result, generated
module, and full compiler diagnostic output.

The independent scope and reconstruction probe is run with:
\begin{lstlisting}
cabal exec -- runghc -package=djex test-church/probe_scopes.hs
\end{lstlisting}
It sends 38 positive and 12 negative type queries through each engine,
covering alpha-equivalence, shadowed binders, nested positive quantifiers,
ambient-variable correlations, higher-kinded application patterns, fresh
polymorphic argument construction, composition and dependent nesting, and
successive quantified provider results exposed after ordinary term arguments.
The
negative cases include direct, indirect, nested, and higher-kinded scope
escape attempts. An absent candidate is recorded only as bounded search
evidence; a timeout is not accepted as a passing negative result. Positive
queries receive 10,000 steps or choices, and every negative query receives
the same smaller budget of 1,000. All retain a three-second wall-clock
safety timeout. Every
returned term is checked by GHC, even if its query was expected to have no
candidate. Five cases provide named polymorphic functions through their
signatures only; their total implementation bodies live in a separate
compiler support module whose constructors are hidden. Every generated
candidate module imports exactly its permitted provider names. One additional
case supplies a unit value; the remaining cases receive no named value
providers.

On the validated build, all 100 queries met their expected outcomes: 76
positive implementations passed GHC, and 24 negative queries returned no
candidate under their search bounds, without errors or timeouts. These
adversarial checks are additional evidence, separate from the 700 Church
implementations and their alias/default wrappers.

The complete Djex validation covers all 19 registered test components: 1,930
ordinary test cases, separately from the 700 Church synthesis cases and 100
scope queries reported above. An aggregate run and a rerun of the repaired
subset close all components successfully. This includes the existing API,
frontend, proof-search, semantic Length, and command-line regression suites,
so the corpus result is accompanied by checks of the established behavior.

Windows validation retains platform-appropriate absolute executable paths and
the corresponding exact policy fingerprints. A short-deadline test also uses
a fake worker that waits without later writes after recording its status
wait. A delayed response could otherwise begin a trace write during process
termination. The test keeps its 200-millisecond query deadline, 2,000-millisecond
shared allowance, required event records, and strict trace parser. These
portability and fixture corrections are separate from the rank-N rules.

From the Leant repository root after building the current executable, run each
engine independently and retain separate reports:
\begin{lstlisting}
python test-church/run_corpus.py --manifest ../Djex/test-church/manifest.json \
  --leant PATH_TO_LEANT_EXECUTABLE --engine djinn --window 1 \
  --steps 4096 --timeout 30 --output test-church/generated-djinn-acceptance
python test-church/run_corpus.py --manifest ../Djex/test-church/manifest.json \
  --leant PATH_TO_LEANT_EXECUTABLE --engine exference --window 1 \
  --steps 4096 --timeout 30 --output test-church/generated-exference-acceptance
\end{lstlisting}
The normal manifest path is \path{lib/Djex/test-church/manifest.json}; the
explicit path above supports development with the separate Djex checkout.
The Lean runner's JSON report and standalone replay output, rather than the
Haskell counts, are the evidence for its target-language results. The separate
\code{--engine both} option runs Leant's combined search mode; success in that
mode does not establish that either individual engine succeeds alone.

The final independent live runs each produce 350 of 350 candidates, and Lean
4.32.0 accepts all 700 exact displayed terms in standalone replay. Every one
of those declarations has an empty axiom inventory. As in the Haskell inventory, each
engine covers 315 provider-free cases, 16 integer-provider cases, and 19 cases
with an explicit element default. The last group remains conditional on that
default in Lean, while Haskell separately checks its documented partial
wrapper at the original signature.

Both runs use a one-candidate window and a configured thirty-second synthesis
timeout. Exference uses 4,096 search steps. The Lean runner's \code{--steps}
option sets \code{synth-steps}, which applies only to Exference; it does not
impose a Djinn choice-point limit. Djinn retains its default unbounded
choice-point budget, \code{synth-budget off}, subject to the shared wall-clock
search deadline and intrinsic finite planning caps. The configured synthesis
timeout is not an end-to-end limit covering process startup, goal serialization,
or the separate standalone kernel replay.
The runs use Leant implementation revision \code{4757569}, with
Djex revision \code{e2eb71e}, on the same unchanged executable with SHA-256
\begin{center}\small\ttfamily
addfac35b9d82955fc871c177b582a8c043475c0171c22cb17977e0e9f5b9869
\end{center}
The live-synthesis reports are
\path{Leant/test-church/generated-djinn-acceptance/results.json} and
\path{Leant/test-church/generated-exference-acceptance/results.json}. They
record source and manifest hashes, executable stability, every candidate,
standalone kernel results, and the axiom inventories. This is fresh synthesis
and exact-term replay for each backend, not a replay-only substitution for a
missing engine run.

The compact fixtures run with:
\begin{lstlisting}
python test-church/run_fixtures.py --leant PATH_TO_LEANT_EXECUTABLE
\end{lstlisting}
Their inventory contains nine representative Church queries, twenty-three broad
rank-N targets repeated in each of the three engine modes, six exact
wide-provider queries, and six named-provider queries with successive
quantified result layers, for 90 queries altogether. The broad targets include genuine
rank-seven continuation encodings with explicit and implicit quantifiers,
composition and dependent construction, and successive quantified result
layers, including ambient open choices and explicit/implicit source and target
quantifier combinations. The eight- and twelve-binder provider tests measure argument width,
which is distinct from type rank. The named layered factories additionally
exercise discovery, source reconstruction, and repeated impredicative choices
at global application sites. These twelve provider queries intentionally
declare opaque primitive premises and require exactly that declared axiom
inventory; the closed Church and broad rank-N fixtures require empty axiom
inventories.

All 90 compact fixture queries pass on that same unchanged executable.
The standalone Lean 4.32.0 replay accepts every exact displayed term. Of these,
78 declarations have empty axiom inventories, and twelve have exactly their
documented supplied-premise inventories. The report is
\path{Leant/test-church/generated-fixtures-bounded-final/results.json}.

The broader ordinary transcript inventory contains 30 files, 263 explicit-type
queries, and two proof-mode \code{:synth} commands, for 265 synthesis commands
in total. Those static counts describe its coverage, separately from any
particular replay result. The Python acceptance harness also checks the order
of interleaved option changes and queries, so a later setting cannot silently
alter an earlier query's test conditions; its 14 harness regression tests pass.
The ordinary Leant Haskell unit suite additionally passes all 558 tests. These
unit and harness checks are distinct from the standalone target-language
replays of the compact fixtures and full corpus.

The full ordinary compatibility run completed on the same unchanged executable
on 5 September 2026 at 01:55:44, UTC$-07{:}00$. It used a configured
600-second synthesis timeout and took 5,515.86 seconds in total. Its original
runner exit code was 1: seventeen fixture outputs matched immediately and
thirteen differed from their recorded baselines. That original result is
preserved; it is not rewritten as an initially passing run.

The thirteen changed fixtures were reviewed query by query: 114 actual
displayed terms across 34 changed queries passed independent Lean 4.32.0
replay. Thirty-one have empty axiom inventories; the other 83 have exactly
their fixture-declared premises. No previously successful query lost its
result. The aggregate record is
\path{Leant/test-church/generated-final-drift-index/results.json}; it excludes
four earlier reviewed baseline updates and their separate 44 term replays.
Only the reviewed term and diagnostic changes were applied to the golden files.

An offline comparison then checked all thirty final golden files against the
retained outputs of that completed live run, covering all 265 synthesis
commands. It used the production Bash runner's normalization functions and
command-substitution semantics; all thirty comparisons pass. This is an
offline comparison of completed live outputs, not a second live synthesis
run. The evidence is retained under
\path{Leant/test-church/generated-goldens-final/}: \path{results.json}
records the original exit status, timing, and executable identity;
\path{final-comparison.json} records the final per-fixture comparisons;
\path{runner.log} and \path{transcripts/} retain the live output. This ordinary
compatibility closure is separate from the independently checked 700-case
Lean corpus and 90 compact fixture results above.

\clearpage
\section{Limits, diagnostics, and future extensions}

The changes remove concrete structural barriers; they do not remove the need
for bounded search. The following limits remain explicit:
\begin{itemize}
  \item The directed Djinn route handles context-free leading schemes.
    Existing constrained-provider and dictionary rules retain their own
    eligibility conditions. Structural context equality in Exference is not
    constraint solving.
  \item The directed matcher uses supplied query structure and a finite
    vocabulary for unconstrained parameters. It does not enumerate every
    possible intermediate polymorphic type.
  \item Historical Djinn instance families, retained-axiom counts, tuple
    attempts, Exference queue limits, and user time/step budgets can still make
    a search inconclusive.
  \item Optional Lean provider-assignment discovery has a shared per-provider
    attempt allowance and a local heartbeat budget. A provider may therefore
    remain available without every useful exact assignment being discovered;
    missing optional evidence is not a non-inhabitation result.
  \item General higher-rank subsumption, polymorphic-let inference, and
    arbitrary higher-order type unification are not claimed as complete
    algorithms. Scope-safe quantified structural equality is a narrower,
    auditable operation.
  \item Partial source signatures require an elaborated expected type. The
    corpus uses GHC to resolve them; a bare underscore without program context
    is not a complete synthesis specification.
  \item Lean validates universe-correct counterparts. The translation does
    not identify \code{Type u} with \code{Type (u + 1)}, assert cumulative
    universes, or replace computational types by propositions.
  \item Corpus success means matching type signatures under the documented
    assumptions. It says nothing by itself about runtime cost, termination of
    imported providers, or equivalence to the source algorithms.
\end{itemize}

A useful future extension should add a general rule, preserve the scope and
evidence invariants, include a failing practical example, and replay emitted
terms in the actual target language. Expanding a numeric cutoff alone is not
a substitute for understanding which structural information a query already
provides.

\clearpage
\appendix
\section{Implementation and artifact map}

\begingroup\small
\begin{longtable}{>{\raggedright\arraybackslash}p{0.37\textwidth}>{\raggedright\arraybackslash}p{0.56\textwidth}}
\toprule
Component & Responsibility \\
\midrule
\endhead
\path{Djinn/Internal/}\newline\path{DirectedInstantiation.hs} &
Structural demand matching, correlated leading-binder selection, nested
scope correspondence, and unobserved-choice scheduling. \\
\path{Djinn/Internal/}\newline\path{Instantiation.hs} &
Scheme eligibility, vocabulary scope checks, bridge construction,
deduplication, and operational family bounds. \\
\path{Djinn/Internal/TypeFormula.hs}\newline
\path{Djinn/Internal/Environment.hs} &
Coherent transport-versus-introduction views and query-aware formulas for
real loaded providers. \\
\path{Exference/Core/Unify.hs} &
Structural quantified equations, private rigid binder identities,
capture-safe polytype substitution, and occurrence/escape checks. \\
\path{Synthesis/Generated.hs} &
Closed and partially specified visible type-argument syntax, lexical binder
identities, and matching of emission patterns against selected types. \\
\path{Exference/Core/Internal/}\newline\path{ExpressionCheck.hs} &
Independent checking, completed occurrence evidence, late instantiation
reconstruction, and rechecking of the rewritten expression. \\
\path{Exference/Core/Internal/}\newline\path{Exference.hs}\newline
\path{Polytype.hs} &
Result-directed selection across successive quantified result layers, with
actual argument obligations retained by ordinary search. \\
\path{Exference/Core/Expression.hs} &
Capture-avoiding normalization of let aliases used at visible
type-application boundaries. \\
\path{test-church/extract_corpus.py} &
Source inventory, GHC hole resolution, alias expansion, contextual defaults,
and artifact freshness. \\
\path{test-church/manifest.json} &
Full provenance, source and resolved signatures, expanded ASTs, scope, and
case classification. \\
\path{test-church/Main.hs} &
Public-API synthesis through both Haskell backends and independent GHC
checking against expanded and alias-preserving types. \\
\path{test-church/probe_scopes.hs} &
Independent alpha-renaming, scope, correlation, higher-kinded application,
and nested-source-reconstruction queries with GHC replay. \\
Leant:\newline\path{src/Leant/Synth/Render.hs}\newline
\path{src/Leant/Synth/Engine.hs} &
Expected-type-guided Lean reconstruction, retained let schemes, constructor
fields, lambda boundaries, and the checked search fallback. \\
Leant:\newline\path{src/Leant/Synth/Fragment.hs} &
Type reification, exact provider metadata, and separate transport resource
bounds. \\
Leant:\newline\path{src/Leant/Backend.hs} &
Owned backend processes, bounded response reading, request deadlines, and
cleanup. \\
Leant:\newline\path{test-church/run_corpus.py} &
Universe-aware Lean goals, explicit provider policy, exact candidate capture,
and standalone kernel replay. \\
Leant:\newline\path{test-church/UniverseProbe.lean} &
Small universe and polymorphic-argument examples checked by Lean. \\
Leant:\newline\path{test-church/run_fixtures.py} &
Exact displayed-term replay for compact Church, broad rank-N, live
eight- and twelve-argument providers, and named layered factories. \\
\bottomrule
\end{longtable}
\endgroup

\begin{thebibliography}{9}
\bibitem{quicklook}
Alejandro Serrano, Jurriaan Hage, Simon Peyton Jones, and Dimitrios Vytiniotis.
\emph{A Quick Look at Impredicativity}.
Proceedings of the ACM on Programming Languages 4 (ICFP), article 89, 2020.
\href{https://doi.org/10.1145/3408971}{doi:10.1145/3408971}.
\href{https://simon.peytonjones.org/assets/pdfs/impredicativity.pdf}{Author-hosted paper}.

\bibitem{inhabitation}
Andrej Dudenhefner and Jakob Rehof.
\emph{A Simpler Undecidability Proof for System F Inhabitation}.
TYPES 2018, LIPIcs 130, 2:1--2:11, published 2019.
\href{https://drops.dagstuhl.de/entities/document/10.4230/LIPIcs.TYPES.2018.2}
{doi:10.4230/LIPIcs.TYPES.2018.2}.

\bibitem{wells}
J.~B. Wells.
\emph{Typability and Type Checking in System F Are Equivalent and Undecidable}.
Annals of Pure and Applied Logic 98(1--3), 111--156, 1999.
\href{https://doi.org/10.1016/S0168-0072(98)00047-5}
{doi:10.1016/S0168-0072(98)00047-5}.
\href{https://people.eecs.berkeley.edu/~necula/Papers/WellsSystemFUndecidable.pdf}
{Earlier conference paper}.

\bibitem{ghc}
GHC contributors.
\emph{GHC 9.12.4 User's Guide, section 6.4.21: Impredicative polymorphism}.
\href{https://downloads.haskell.org/ghc/9.12.4/docs/users_guide/exts/impredicative_types.html}
{Versioned GHC reference chapter}.

\bibitem{leanuniverses}
Lean contributors.
\emph{The Lean Language Reference: Universes}.
\href{https://lean-lang.org/doc/reference/latest/The-Type-System/Universes/}
{Lean reference chapter}.
\end{thebibliography}
\end{document}
